% intro.tex
\clearpage
{
  \fontsize{11}{13.5}\selectfont
  \setlength{\parindent}{1.5em}
\section*{Foreword}

There is another significant variant of the text of Esther called the Alpha Text (or AT) that is found four medieval manuscripts and is substantially different from both the Hebrew Masoretic Text and the Greek Septuagint versions. This has been published (in the original Greek) in Hanhart's {\textit Septuaginta:
Vetus Testamentum Graecum Auctoritate Academiae Scientiarum Gottingensis editum VIII.3: Esther} [Göttingen:
Vandenhoeck \& Ruprecht, 1966]. Until now the only readily available source for English speakers to access this text has been the excellent New English Translation of the Septuagint (NETS) Esther. This can be found at https://ccat.sas.upenn.edu/nets/edition/17-esther-nets.pdf. Although it is readily accessible and liberally licensed, the NETS Esther is protected by copyright meaning it cannot be freely redistributed or used to make derivative works without jumping through the hoops necessary to obtain copyright permission. For many independent scholars that is prohibitive.

I have produced (with the assistance of generative AI) a new translation of the Alpha text from Hanhart's Greek text (as obtained from a Logos Bible software module) under a Creative Commons Attribution licence. You are essentially free to do whatever you like with this text. (Please let me know if you intend to create a commercial product from it, but no additional permission need be sought.)

Discussions of the historical background of these three text traditions are beyond the scope of this work. Interested readers may wish to consult M.V. Fox. 1990. “The Alpha Text of the Greek Esther.” Textus, 15, Pp. 27-54.

This document has as a companion a parallel edition that presents these three textual traditions side by side, using the World English Bible (WEB) translation for the Masoretic Text and Septuagint traditions alongside this translation.

\subsection*{Chapter and verse divisions}

Verse numbering of the Greek version of Esther in English Bibles usually follows the ordering Jerome used in the Vulgate. This provides the Hebrew version of Esther as found in the MT first, with all the additional material from the Greek versions as additional chapters at the end. These additions are labelled A to F. In this edition, the chapter and verse numbering is as it appears in the Göttingen edition, but with indicators showing where the text corresponds with the chapter divisions and additions in the Septuagint. From about chapter 7 the Alpha Text diverges from the other sources, and the Göttingen edition provides no further chapter divisions, simply numbering verses sequentially from the start of chapter 7. At the beginning of many paragraphs a second verse reference appears in parentheses to show which verse in the usual English version of Esther (from the MT) corresponds to this verse in the AT. This is absent for some paragraphs as some verses in the AT, especially in the additional chapters, have no corresponding passage in the MT.

\subsection*{Spelling}
The World English Bible is available in a British English edition as well as its standard American English edition. Both versions are available for the MT and LXX columns of this document. The third column translated from the Alpha Text also comes in US and British versions, with the only difference being scripted spelling adjustments. This is the \ifbritish{British}\else{American}\fi\ English version.

\bigskip
\noindent Glen Prideaux\\
March 2026
}